%! AP Statistics — Unit 9, Lesson 9.5 — Student Packet Cover Sheet
\documentclass[10pt]{article}
\usepackage{apstats-article}
\usepackage{apstats-boxes}
\usepackage{tikz}

\begin{document}

\begin{tikzpicture}[remember picture, overlay]
  \fill[navy] (current page.north west) rectangle
    ([yshift=-2.2cm]current page.north east);
  \node[anchor=west, white, font=\large\bfseries] at
    ([xshift=1cm, yshift=-0.7cm]current page.north west)
    {\CourseName: \SchoolYear};
  \node[anchor=west, white, font=\normalsize] at
    ([xshift=1cm, yshift=-1.3cm]current page.north west)
    {Unit 9: Inference for Quantitative Data: Slopes};
  \node[anchor=west, white, font=\normalsize\bfseries] at
    ([xshift=1cm, yshift=-1.9cm]current page.north west)
    {Lesson 9.5 \quad Carrying Out a Test for the Slope of a Regression Model};
\end{tikzpicture}

\vspace{2.4cm}
\namedateperiod

\begin{learningtargetbox}
\begin{itemize}
    \item I can calculate the $t$-test statistic for the slope of a regression model
          using $t = b / SE_b$.
    \item I can interpret the $p$-value of a significance test for slope in context.
    \item I can carry out all four steps of a significance test for slope and state
          a conclusion.
    \item I can read a regression output table and identify $b$, $SE_b$, $t$, and
          the $p$-value.
\end{itemize}
\end{learningtargetbox}

\begin{tocbox}
\renewcommand{\arraystretch}{1.5}
\begin{tabularx}{\linewidth}{clXc}
    \textbf{\#} & \textbf{Component} & \textbf{Description} & \textbf{Score} \\
    \hline
    1 & Warm-Up      & Spiral review: $t$-statistic, degrees of freedom, decision rule & \hfill/6 \\
    2 & Guided Notes & Reading output tables, 4-step significance test for slope & \hfill/12 \\
    3 & Activity     & Two regression contexts: exam scores and body mass index & \hfill/18 \\
    4 & Exit Ticket  & Independent: hypotheses, $p$-value interpretation, conclusion & \hfill/6 \\
    5 & Homework     & Practice carrying out full 4-step tests for slope & \hfill/12 \\
    \hline
       & \multicolumn{2}{l}{\textbf{Total}} & \hfill/54 \\
\end{tabularx}
\end{tocbox}

\end{document}
